%% ============================================================
%%  Capítulo 4 — Desenho e Especificação
%% ============================================================
\chapter{Desenho e Especificação}
\label{chap:desenho}

\section{Arquitetura do Sistema}

A arquitetura segue o padrão \textbf{Resource Server + Identity Provider}, com separação clara de responsabilidades entre os três serviços orquestrados por Docker Compose.

\begin{figure}[h!]
  \centering
  \includegraphics[width=0.95\textwidth]{arquitetura_iam.png}
  \caption{Arquitetura do projeto: FastAPI, Keycloak 26 e PostgreSQL 16 sob Docker Compose.}
  \label{fig:arquitetura}
\end{figure}

O fluxo principal, ilustrado na Figura~\ref{fig:arquitetura}, é: (1)~o cliente autentica no Keycloak e recebe um \gls{jwt} assinado com RS256; (2)~apresenta o token à FastAPI; (3)~a FastAPI valida a assinatura localmente via \gls{jwks} cacheado; (4)~aplica \gls{rbac}; (5)~responde ao cliente.

\begin{table}[h!]
\centering
\caption{Componentes da arquitetura, portos e funções}
\label{tab:componentes}
\begin{tabular}{|l|c|l|}
\hline
\textbf{Componente} & \textbf{Porto} & \textbf{Função principal} \\ \hline
Keycloak~26 & 8081 & \gls{idp}: emissão de \gls{jwt}, realm, \gls{rbac}, \gls{mfa}, auditoria \\ \hline
FastAPI (Python~3.12) & 8000 & \textit{Resource server}: validação \gls{jwt}, endpoints protegidos, dashboard \\ \hline
PostgreSQL~16 & 5432 & Base de dados persistente do Keycloak \\ \hline
JML Scripts (\gls{cli}) & — & Automação do ciclo de vida via Admin \gls{api} \\ \hline
\end{tabular}
\end{table}

\subsection{Reprodutibilidade da Infraestrutura}

Toda a configuração é declarativa e versionada em git:

\begin{itemize}
  \item \textbf{\texttt{docker-compose.yml}} — define os três serviços, redes internas, volumes e \textit{health checks} sequenciais;
  \item \textbf{\texttt{realm-export.json}} — exportação completa do \textit{realm} \texttt{iam-tp}: papéis, grupos, utilizadores, cliente \gls{oidc} e fluxos de \gls{mfa};
  \item \textbf{\texttt{.env}} — credenciais sensíveis (não versionado), seguindo o princípio de não expor \textit{secrets} no código.
\end{itemize}

\begin{lstlisting}[caption={Arranque completo da infraestrutura (primeira execução: 2--3 min)},label={lst:docker}]
cp .env.example .env          # criar ficheiro de secrets
docker compose up --build     # iniciar todos os servicos
curl http://localhost:8000/health
# Resposta esperada: {"status":"ok"}
\end{lstlisting}

\section{Autenticação OIDC e Validação JWT}
\label{sec:autenticacao}

\subsection{Fluxo de Autenticação}

A autenticação é completamente delegada no Keycloak. A FastAPI actua como \textit{Resource Server}, nunca vendo as credenciais do utilizador. Para a demonstração usa-se o fluxo \gls{ropc} via \texttt{/auth/token}; o realm está também configurado com \textit{Authorization Code + PKCE} para fluxos \textit{browser} em produção.

\subsection{Validação JWT RS256/JWKS}

A validação é feita \textbf{localmente} na FastAPI, sem chamadas ao Keycloak em cada pedido:

\begin{enumerate}
  \item No primeiro pedido autenticado, a FastAPI descarrega a \gls{jwks} publicada pelo \textit{realm};
  \item A chave pública RS256 é cacheada em memória (\texttt{\_jwks\_cache} em \texttt{app/auth.py});
  \item Cada token é verificado localmente: assinatura criptográfica, \texttt{issuer}, \texttt{audience} e expiração;
  \item As \textit{claims} relevantes são extraídas do \textit{payload} para autorização.
\end{enumerate}

\begin{lstlisting}[language=Python, caption={Validação JWT RS256 em app/auth.py},label={lst:jwt}]
async def get_current_user(
    token: str = Depends(oauth2_scheme)
) -> dict:
    jwks = await _get_jwks()       # cache em memoria
    payload = jwt.decode(
        token,
        jwks,
        algorithms=["RS256"],
        audience=settings.oidc_client_id,
        options={"verify_exp": True},
    )
    return payload                 # claims disponiveis nos handlers
\end{lstlisting}

\subsection{Claims Utilizadas pela Aplicação}

\begin{table}[h!]
\centering
\caption{Claims JWT extraídas e utilizadas pela aplicação}
\label{tab:claims}
\begin{tabular}{|l|l|l|}
\hline
\textbf{Claim} & \textbf{Descrição} & \textbf{Uso} \\ \hline
\texttt{sub} & UUID único do utilizador (Keycloak) & Auditoria, perfil \\ \hline
\texttt{email} & Endereço de correio eletrónico & Exibição no perfil \\ \hline
\texttt{preferred\_username} & Nome de utilizador de login & Identificação visual \\ \hline
\texttt{realm\_access.roles} & Lista de papéis do \textit{realm} & Base do \gls{rbac} \\ \hline
\texttt{amr} / \texttt{acr} & Método de autenticação (\gls{mfa}) & Verificação \gls{mfa} \\ \hline
\texttt{exp} & \textit{Timestamp} de expiração & Validade do token \\ \hline
\end{tabular}
\end{table}

\section{Autorização RBAC}
\label{sec:autorizacao}

\subsection{Modelo de Papéis e Grupos}

Foram definidos três papéis de \textit{realm} no Keycloak, com grupos correspondentes que permitem atribuição automática de papéis via \textit{group membership}:

\begin{itemize}
  \item \textbf{\texttt{admin}} — Grupo \texttt{/Admins}: acesso total, incluindo \gls{jml}, auditoria e zona \gls{mfa}. \gls{mfa} \gls{totp} obrigatória;
  \item \textbf{\texttt{colaborador}} — Grupo \texttt{/Colaboradores}: acesso a recursos e dados internos;
  \item \textbf{\texttt{visitante}} — Grupo \texttt{/Visitantes}: acesso apenas a recursos públicos e perfil próprio.
\end{itemize}

\begin{figure}[h!]
  \centering
  \includegraphics[width=0.95\textwidth]{matriz_rbac.png}
  \caption{Matriz RBAC: mapeamento de endpoints para papéis com resultados esperados (200/403).}
  \label{fig:rbac}
\end{figure}

\subsection{Implementação das Dependências de Autorização}

\begin{table}[h!]
\centering
\caption{Controlo de acesso por endpoint}
\label{tab:endpoints}
\rowcolors{2}{gray!15}{white}
\begin{tabular}{|l|l|c|c|c|}
\hline
\textbf{Endpoint} & \textbf{Proteção} & \textbf{admin} & \textbf{colab.} & \textbf{visit.} \\ \hline
\texttt{GET /public} & Nenhuma & 200 & 200 & 200 \\ \hline
\texttt{GET /me} & JWT válido & 200 & 200 & 200 \\ \hline
\texttt{GET /colaborador/data} & colaborador\textbar{}admin & 200 & 200 & 403 \\ \hline
\texttt{GET /admin/users} & admin & 200 & 403 & 403 \\ \hline
\texttt{GET /admin/audit} & admin & 200 & 403 & 403 \\ \hline
\texttt{GET /admin/mfa-area} & admin + \gls{mfa} & 200* & 403 & 403 \\ \hline
\multicolumn{5}{l}{\footnotesize * apenas se o token incluir \texttt{amr=otp}} \\
\end{tabular}
\end{table}

\begin{lstlisting}[language=Python, caption={Dependências require\_role() e require\_mfa() em auth.py},label={lst:rbac}]
def require_role(*roles: str):
    async def checker(user=Depends(get_current_user)):
        user_roles = user.get("realm_access", {}).get("roles", [])
        if not any(r in user_roles for r in roles):
            raise HTTPException(status_code=403)
        return user
    return checker

def require_mfa():
    async def checker(user=Depends(get_current_user)):
        amr = user.get("amr", [])
        if not any(m in amr for m in ["otp", "mfa", "2fa"]):
            raise HTTPException(status_code=403)
        return user
    return checker
\end{lstlisting}

\section{Autenticação Multifator (MFA/TOTP)}
\label{sec:mfa}

A \gls{mfa} foi implementada com \gls{totp} \cite{RFC6238}, compatível com Google Authenticator, Microsoft Authenticator e Authy.

\subsection{Configuração no Keycloak}

\begin{enumerate}
  \item O fluxo de autenticação \textit{browser} inclui a execução \texttt{auth-otp-form} (obrigatória para o papel \texttt{admin});
  \item O utilizador \texttt{admin.user} tem a \textit{Required Action} \texttt{CONFIGURE\_TOTP} ativa;
  \item No primeiro acesso via \textit{browser}, o Keycloak força a leitura do QR Code e a configuração da aplicação TOTP;
  \item Após configuração, o \gls{jwt} passa a incluir \texttt{"amr": ["otp"]} como prova de segundo fator.
\end{enumerate}

\subsection{Validação na API}

O endpoint \texttt{/admin/mfa-area} exige cumulativamente o papel \texttt{admin} e a \textit{claim} \texttt{amr=otp}. Um token de \texttt{admin.user} obtido via \gls{ropc} (sem \gls{mfa}) retorna \textbf{403}; um token obtido via \textit{browser flow} com \gls{totp} retorna \textbf{200}. Esta separação demonstra que a \gls{api} valida a prova de \gls{mfa} de forma independente do Keycloak.

\section{Ciclo de Vida JML}
\label{sec:jml}

\begin{figure}[h!]
  \centering
  \includegraphics[width=0.95\textwidth]{fluxo_jml.png}
  \caption{Fluxos JML: criação, transição de papel/grupo, desativação e revogação de sessões.}
  \label{fig:jml}
\end{figure}

Os três fluxos foram implementados sobre a Keycloak Admin \gls{rest} \gls{api}, garantindo que toda a gestão de identidades é centralizada no \gls{idp}. Cada fluxo está disponível como script \gls{cli} (\texttt{jml/joiner.py}, \texttt{mover.py}, \texttt{leaver.py}) e como endpoint \gls{rest} (\texttt{POST /admin/jml/\{joiner|mover|leaver\}}).

\begin{table}[h!]
\centering
\caption{Fluxos JML: ações e efeito de segurança imediato}
\label{tab:jml}
\begin{tabularx}{\textwidth}{|l|X|l|}
\hline
\textbf{Fluxo} & \textbf{Ações} & \textbf{Efeito imediato} \\ \hline
Joiner & Cria utilizador, atribui papel e grupo inicial & Acesso inicial controlado \\ \hline
Mover & Troca papel/grupo, \textbf{revoga sessões ativas} & Novo acesso após reautenticação \\ \hline
Leaver & Desativa conta (\texttt{enabled:false}), remove papéis, \textbf{revoga sessões} & Bloqueio imediato de acesso \\ \hline
\end{tabularx}
\end{table}

\subsection{Decisão Técnica: Revogação Explícita de Sessões}

Sem revogação explícita, um utilizador movido ou desativado manteria o acesso durante o tempo de vida do \gls{jwt} (por omissão, 300 segundos). A chamada \texttt{DELETE /admin/realms/iam-tp/users/\{id\}/sessions} força a reautenticação imediata, eliminando esta janela de acesso indevido.

\begin{lstlisting}[language=Python, caption={Revogação de sessões no mover.py},label={lst:revoke}]
# Após trocar papel e grupo:
admin_token = get_admin_token()
user_id = get_user_id(username, admin_token)
revoke_url = f"{KEYCLOAK_URL}/admin/realms/{REALM}/users/{user_id}/sessions"
httpx.delete(revoke_url, headers={"Authorization": f"Bearer {admin_token}"})
\end{lstlisting}

\section{Auditoria de Eventos}
\label{sec:auditoria}

O Keycloak regista automaticamente dois tipos de eventos com persistência em PostgreSQL:

\begin{itemize}
  \item \textbf{Eventos de utilizador}: login, falha de login, logout, \textit{refresh} de token, configuração de \gls{totp} e falhas de \gls{otp};
  \item \textbf{Eventos de administração}: criação/remoção de utilizadores, alteração de papéis e grupos, atualização de clientes.
\end{itemize}

O endpoint \texttt{GET /admin/audit} (protegido por papel \texttt{admin}) agrega ambos os tipos numa única resposta e suporta filtragem por tipo de evento.

\subsection{Excerto de Evidência de Auditoria}

A Listagem~\ref{lst:audit_evento} apresenta um excerto real de evento de login registado pelo Keycloak, tal como retornado pelo endpoint \texttt{/admin/audit}:

\begin{lstlisting}[caption={Excerto de evento de auditoria — login bem-sucedido},label={lst:audit_evento}]
{
  "id": "b3a7c1f2-8d4e-4a2b-9f1c-3e5d7a8b2c4d",
  "time": 1748430600000,
  "type": "LOGIN",
  "realmId": "iam-tp",
  "clientId": "fastapi-client",
  "userId": "a1b2c3d4-e5f6-7890-abcd-ef1234567890",
  "sessionId": "c9d8e7f6-a5b4-3c2d-1e0f-9a8b7c6d5e4f",
  "ipAddress": "172.18.0.1",
  "details": {
    "auth_method": "openid-connect",
    "auth_type": "code",
    "response_type": "code",
    "redirect_uri": "http://localhost:8000/auth/callback",
    "consent": "no_consent_required",
    "code_id": "...",
    "username": "colaborador.user",
    "response_mode": "query"
  }
}
\end{lstlisting}

E um excerto de evento administrativo \gls{jml} (criação de utilizador via Joiner):

\begin{lstlisting}[caption={Excerto de evento de administração — criação de utilizador (Joiner)},label={lst:audit_admin}]
{
  "id": "f4e3d2c1-b0a9-8765-4321-fedcba987654",
  "time": 1748431200000,
  "operationType": "CREATE",
  "resourceType": "USER",
  "realmId": "iam-tp",
  "representation": "{\"username\":\"demo.user\",\"enabled\":true}",
  "resourcePath": "users/f1e2d3c4-a5b6-7890-abcd-ef0123456789",
  "authDetails": {
    "realmId": "iam-tp",
    "clientId": "admin-cli",
    "userId": "admin-service-account",
    "ipAddress": "172.18.0.1"
  }
}
\end{lstlisting}

A evidência consultável completa está disponível no dashboard interativo (aba \textbf{Auditoria}), com filtragem por tipo de evento (\texttt{LOGIN}, \texttt{LOGIN\_ERROR}, \texttt{UPDATE\_TOTP}, etc.) e resumo estatístico por categoria.

\section{Segurança Operacional}

\begin{table}[h!]
\centering
\caption{Riscos identificados e mitigações implementadas}
\label{tab:riscos}
\rowcolors{2}{gray!15}{white}
\begin{tabularx}{\textwidth}{|X|X|}
\hline
\textbf{Risco} & \textbf{Mitigação} \\ \hline
Token inválido ou forjado & Validação RS256 via \gls{jwks}, verificação de \texttt{issuer} e \texttt{audience} \\ \hline
Acesso indevido a recursos & \gls{rbac} no \textit{backend} com \texttt{require\_role()} por endpoint \\ \hline
Administrador sem segundo fator & \gls{totp} obrigatório e validação de \texttt{amr} na \gls{api} \\ \hline
Sessões ativas após mudança de papel & Revogação explícita (\texttt{DELETE /users/\{id\}/sessions}) no Mover e Leaver \\ \hline
Credenciais expostas no código & Variáveis de ambiente via \texttt{.env} (não versionado) \\ \hline
Falta de rastreabilidade & Eventos Keycloak (PostgreSQL) e endpoint \texttt{/admin/audit} protegido \\ \hline
\end{tabularx}
\end{table}
